\section{Mục tiêu của đề tài}
Mục tiêu tổng quát của khóa luận là thiết kế và phát triển một hệ thống IoT có khả năng giám sát chất lượng không khí trong căn hộ, theo dõi điều kiện chăm sóc cây trầu bà và điều khiển một số thiết bị hỗ trợ cải thiện môi trường sống. Định hướng của đề tài là chứng minh tính khả thi của mô hình kết hợp IAQ, trầu bà/thủy canh quy mô nhỏ và điều khiển thiết bị theo thời gian thực.

Các mục tiêu cụ thể gồm:
\begin{itemize}
    \item Thiết kế node cảm biến sử dụng Arduino Pro Mini để đọc TDS, pH, nhiệt độ, độ ẩm, CO2 và ánh sáng.
    \item Xây dựng kênh truyền LoRa giữa node cảm biến và gateway ESP32.
    \item Phát triển gateway ESP32 để nhận dữ liệu LoRa, kiểm tra gói tin và gửi telemetry lên ThingsBoard.
    \item Thiết kế ứng dụng Flutter kết nối ThingsBoard để hiển thị dữ liệu, trạng thái cảnh báo, lịch sử thông số và cấu hình ngưỡng vận hành.
    \item Xây dựng chức năng điều khiển relay bật/tắt quạt, đèn và điều khiển tăng giảm nhiệt độ điều hòa qua IR tại node.
    \item Đề xuất các luật điều khiển dựa trên ngưỡng CO2, ánh sáng, nhiệt độ, độ ẩm, TDS và pH nhằm hỗ trợ cải thiện IAQ và duy trì điều kiện phù hợp cho trầu bà.
    \item Đánh giá khả năng hoạt động của hệ thống trong kịch bản căn hộ chung cư.
\end{itemize}
